\begin{problem}[第31届全国中学生物理竞赛复赛][失重液滴脉动]{17}
2013年6月20日，“神舟十号”女航天员王亚平在“天宫一号”目标飞行器里成功进行了我国首次太空授课。授课中的一个实验展示了失重状态下液滴的表面张力引起的效应。视频中可发现漂浮的液滴处于周期性的“脉动”中（平时在地球表面附近，重力的存在会导致液滴下降太快，以至于很难观察到液滴的这种“脉动”现象）。假设液滴处于完全失重状态，液滴的上述“脉动”可视为液滴形状的周期性的微小变化（振动）。
    \begin{subproblem}
        该液滴处于平衡状态时的形状是什么？
    \end{subproblem}
    \begin{subproblem}
        决定该液滴振动频率 $f$ 的主要物理量是什么？
    \end{subproblem}
    \begin{subproblem}
        按后面括号中提示的方法导出液滴振动频率与上述物理量的关系式。（提示：例如，若认为 $a, b, c$ 是决定该液滴振动频率的相互独立的主要物理量，可将液滴振动频率 $f$ 与 $a, b, c$ 的关系式表示为 $f \propto a^{\alpha} b^{\beta} c^{\gamma}$ ，其中指数 $\alpha, \beta, \gamma$ 是相应的待定常数。）
    \end{subproblem}
\end{problem}